% Options for packages loaded elsewhere
\PassOptionsToPackage{unicode}{hyperref}
\PassOptionsToPackage{hyphens}{url}
\documentclass[
]{article}
\usepackage{xcolor}
\usepackage[margin=1in]{geometry}
\usepackage{amsmath,amssymb}
\setcounter{secnumdepth}{-\maxdimen} % remove section numbering
\usepackage{iftex}
\ifPDFTeX
  \usepackage[T1]{fontenc}
  \usepackage[utf8]{inputenc}
  \usepackage{textcomp} % provide euro and other symbols
\else % if luatex or xetex
  \usepackage{unicode-math} % this also loads fontspec
  \defaultfontfeatures{Scale=MatchLowercase}
  \defaultfontfeatures[\rmfamily]{Ligatures=TeX,Scale=1}
\fi
\usepackage{lmodern}
\ifPDFTeX\else
  % xetex/luatex font selection
\fi
% Use upquote if available, for straight quotes in verbatim environments
\IfFileExists{upquote.sty}{\usepackage{upquote}}{}
\IfFileExists{microtype.sty}{% use microtype if available
  \usepackage[]{microtype}
  \UseMicrotypeSet[protrusion]{basicmath} % disable protrusion for tt fonts
}{}
\makeatletter
\@ifundefined{KOMAClassName}{% if non-KOMA class
  \IfFileExists{parskip.sty}{%
    \usepackage{parskip}
  }{% else
    \setlength{\parindent}{0pt}
    \setlength{\parskip}{6pt plus 2pt minus 1pt}}
}{% if KOMA class
  \KOMAoptions{parskip=half}}
\makeatother
\usepackage{graphicx}
\makeatletter
\newsavebox\pandoc@box
\newcommand*\pandocbounded[1]{% scales image to fit in text height/width
  \sbox\pandoc@box{#1}%
  \Gscale@div\@tempa{\textheight}{\dimexpr\ht\pandoc@box+\dp\pandoc@box\relax}%
  \Gscale@div\@tempb{\linewidth}{\wd\pandoc@box}%
  \ifdim\@tempb\p@<\@tempa\p@\let\@tempa\@tempb\fi% select the smaller of both
  \ifdim\@tempa\p@<\p@\scalebox{\@tempa}{\usebox\pandoc@box}%
  \else\usebox{\pandoc@box}%
  \fi%
}
% Set default figure placement to htbp
\def\fps@figure{htbp}
\makeatother
\setlength{\emergencystretch}{3em} % prevent overfull lines
\providecommand{\tightlist}{%
  \setlength{\itemsep}{0pt}\setlength{\parskip}{0pt}}
\usepackage{bookmark}
\IfFileExists{xurl.sty}{\usepackage{xurl}}{} % add URL line breaks if available
\urlstyle{same}
\hypersetup{
  hidelinks,
  pdfcreator={LaTeX via pandoc}}

\author{}
\date{\vspace{-2.5em}}

\begin{document}

\section{Sistema de Inventario y
Facturación}\label{sistema-de-inventario-y-facturaciuxf3n}

\begin{center}\rule{0.5\linewidth}{0.5pt}\end{center}

\subsection{Información General}\label{informaciuxf3n-general}

\begin{itemize}
\item
  \textbf{Curso:} Base de Datos I\\
\item
  \textbf{Motor de Base de Datos:} Oracle data base (DBeaver versión
  26.2.0)
\item
  \textbf{Alumna:} Sol Arely Martínez García\\
\item
  \textbf{Carnet:} 7590-24-4829
\item ~
  \subsection{\texorpdfstring{\textbf{Fecha:}
  31/10/2026}{Fecha: 31/10/2026}}\label{fecha-31102026}
\end{itemize}

\subsection{Descripción Breve del Sistema y
Alcance}\label{descripciuxf3n-breve-del-sistema-y-alcance}

\subsubsection{Descripción}\label{descripciuxf3n}

Este proyecto consiste en el diseño e implementación de una base de
datos relacional para gestionar el control de inventario y el proceso de
facturación en un negocio comercial (enfocado en una ferretería). El
sistema permite administrar los datos principales de la tienda
(productos, categorías, proveedores y clientes) y llevar un registro
detallado de cada venta que se realiza.

\subsubsection{Alcance del Sistema}\label{alcance-del-sistema}

\begin{itemize}
\tightlist
\item
  \textbf{Gestión de Catálogos:} Registro y consulta de categorías de
  productos, proveedores y clientes.
\item
  \textbf{Control de Inventario:} Manejo de existencias (stock actual),
  precios de compra y venta, vinculados con sus respectivas categorías y
  proveedores.
\item
  \textbf{Proceso de Facturación:} Registro completo de ventas con
  encabezado (datos de la factura y cliente) y detalle línea por línea
  (productos comprados, cantidades y subtotales).
\item
  \textbf{Automatización Integrada:} Actualización automática del stock
  disponible y cálculo de totales mediante la lógica implementada en la
  base de datos
\item
  \textbf{Reportes y Consultas:} Generación de reportes clave para la
  toma de decisiones, tales como productos con bajo stock, historial de
  ventas por cliente, productos más vendidos e inventario valorizado por
  proveedor.
\end{itemize}

\begin{center}\rule{0.5\linewidth}{0.5pt}\end{center}

\subsection{Instrucciones Rápidas de
Ejecución}\label{instrucciones-ruxe1pidas-de-ejecuciuxf3n}

Para montar la base de datos desde cero en el entorno local con DBeaver
o cliente SQL preferido, seguirestos pasos ejecutando los scripts
ubicados en la carpeta \texttt{sql/} en el orden exacto indicado:

\begin{enumerate}
\def\labelenumi{\arabic{enumi}.}
\item
  \textbf{Crear la Base de Datos:}\\
  Ejecuta el script \texttt{sql/01\_ddl\_crear\_base\_datos.sql} para
  preparar el espacio de trabajo.
\item
  \textbf{Crear la Estructura de Tablas:}\\
  Ejecuta el archivo \texttt{sql/02\_ddl\_crear\_tablas.sql}. Este
  script genera todas las tablas necesarias con sus respectivas llaves
  primarias, llaves foráneas y reglas de validación (restricciones).
\item
  \textbf{Cargar Datos de Prueba:}\\
  Ejecuta \texttt{sql/03\_dml\_datos\_prueba.sql} para poblar el sistema
  con información inicial (catálogos de clientes, proveedores,
  categorías, productos con stock y facturas de ejemplo).
\item
  \textbf{Probar Reportes y Consultas:}\\
  Abre y ejecuta \texttt{sql/04\_consultas\_reportes.sql} para verificar
  el correcto funcionamiento del sistema y visualizar los reportes
  generados.
\end{enumerate}

\subsection{Diagrama
Entidad-Relación}\label{diagrama-entidad-relaciuxf3n}

\begin{itemize}
\tightlist
\item
  \texttt{!{[}Diagrama\ Entidad-Relación{]}(img/modelo\_er.png)}
\end{itemize}

\begin{center}\rule{0.5\linewidth}{0.5pt}\end{center}

\subsection{Índice de Documentos y
Scripts}\label{uxedndice-de-documentos-y-scripts}

```text ├── capturas/\\
│ ├── 01\_lista\_tablas.png │ └── 02\_reporte\_bajo\_stock.png ├──
docs/\\
│ ├── 01\_descripcion\_sistema.md │ ├── 02\_diccionario\_datos.md │ ├──
03\_justificacion\_relaciones.md │ ├── 04\_normalizacion.md │ └──
05\_paso\_a\_paso.md ├── img/\\
│ ├── modelo\_er.png │ └── modelo\_relacional.png └── sql/\\
├── 01\_ddl\_crear\_base\_datos.sql ├── 02\_ddl\_crear\_tablas.sql ├──
03\_dml\_datos\_prueba.sql └── 04\_consultas\_reportes.sql

\end{document}
